\noindent HENRY, Marco Antônio Cottorello. NEVES, Murilo de Souza. SILVA JUNIOR, Rodrigo Macedo da.
\textbf{Virtualização, Nuvem e Contêineres}. 2026.
24 folhas.
(Graduação em Ciência da Computação) -- Universidade Estadual de Londrina,
Londrina, 2026.

\vspace{2\baselineskip}

% Resumo
\begin{center}
  \textbf{RESUMO}
\end{center}

\vspace{\baselineskip}

\noindent A consolidação da computação em nuvem deslocou a infraestrutura de tecnologia da informação do servidor dedicado para recursos provisionados sob demanda, sustentados por camadas sucessivas de abstração. Este trabalho apresenta essa pilha de forma integrada, partindo da definição canônica do NIST, com as cinco características essenciais, os três modelos de serviço e os quatro modelos de implantação, e do critério que efetivamente separa IaaS, PaaS e SaaS, que é a posição da fronteira de controle sobre a pilha de software. Em seguida examina duas implementações situadas na camada de abstração e controle de recursos. A primeira é o OpenStack, cuja arquitetura modular e cujos serviços de identidade, imagem, computação e rede convertem servidores físicos em infraestrutura provisionável por interface programável. A segunda é o Docker, cuja engine decomposta, cujas primitivas de isolamento do kernel e cujos modelos de rede local definem o contêiner como processo isolado sobre kernel compartilhado. A partir dessa base, o trabalho analisa a superfície de ataque decorrente do compartilhamento de kernel, com os mecanismos de contenção de privilégios, as vulnerabilidades de imagens e a proteção do tráfego entre contêineres, e discute a eficiência energética de data centers, com a métrica PUE, a consolidação dinâmica de servidores e a gestão térmica. Conclui-se que as escolhas de isolamento e de densidade de consolidação são interdependentes, de modo que segurança e eficiência energética constituem dimensões acopladas de um mesmo projeto de infraestrutura.

\vfill

\noindent\textbf{Palavras-chave:} Computação em Nuvem. Virtualização. Contêineres. OpenStack. Docker. Eficiência Energética.
